\scanpage{207}
% archive sheet number: 206
\textcircled{3}
\diagram{\begin{tikzcd}[ampersand replacement=\&,column sep=large,row sep=large]
X\arrow[d,"{f\ (\text{propre})}"']\arrow[r,phantom,"{\supset}" description,"{\unsure{g|j}}"']\&U\\
S\&\eta\ \deleted{\unsure{hens|gén}}\arrow[l,hook]\arrow[u,"{f_U}"']
\end{tikzcd}}
$X$ régulier, $U=X-\deleted{\unsure{X}}Y$, $Y=\bigcup Y_\alpha\supset X_0$
\[
\left\{\begin{aligned}
Y_\alpha&\quad\text{div. réguliers}\\
Y&\quad\text{div. à croisements normaux.}
\end{aligned}\right.
\]
\margin{\unsure{en} $Y$}
\margin{$X_0$.}
[\uline{NB} faire pour simplifier $Y=X_0$, $U=X_\eta$]
\[
H^{\unsure{*|\bullet}}(U)\Longleftarrow E_2^{pq}=H^p(X_0,R^q g_*(\mathbb{Q}_{\unsure{\ell|l}X})\vert_{X_0})
\]
\[
R^q g_*(\mathbb{Q}_{\ell X})\simeq\mathop{\bigwedge}^{q}\coprod_\alpha(\mathbb{Q}_\ell(-1)_{Y_\alpha})\simeq\coprod_{\alpha_1<\cdots<\alpha_q}\mathbb{Q}_{\ell Y_{\alpha_1\cdots\alpha_q}}(-q)
\]
donc \deleted{\illegible{1}}
\[
\left\{\begin{aligned}
E_2^{p,q}&=\deleted{\unsure{H^p}}\coprod_{\alpha_1<\cdots<\alpha_q}H^p(\deleted{\unsure{X_0}}(Y_{\alpha_1\cdots\alpha_q})_0)(-q)\qquad\text{si }q\geq1\\
E_2^{p0}&=H^p(X_0)\deleted{\unsure{\mathbb{Q}_\ell}}\quad\text{de poids }\leq\deleted{\unsure{q}}p=i\ \deleted{\text{\unsure{au plus}}}
\end{aligned}\right.
\]
\margin{de poids $p+2q=i+q$}
\[
E_r^{pq}\unsure{-|\longrightarrow}E_r^{p+r,q-r+1}
\]
a) $q-r+1\neq0$\quad poids $p+2q$ et $p+r+2(q-r+1)$\quad nulle si $r\neq2$

b) $q-r+1=0$ \deleted{\unsure{po.} $p+2q$} poids $p+2q$ et $\leq p+r$ il faut pour que $d_r^{pq}\neq0$ que l'on ait $p+2q\leq p+r$ i.e. $r\geq2q$; comme $r=q+1$, ceci s'écrit $q+1\geq2q$, i.e. $q\leq1$ donc $r\leq2$.

Donc $\underline{d_r=0}$ si $r\unsure{\geq|\gg}3$.

\deleted{En \unsure{diminuant|dimension}}

$H^i$ filtration \unsure{dont} les termes, en croissant, sont
\[
\begin{aligned}
E_2^{i0}&\quad\text{de poids }\deleted{\unsure{\leq}}i\\
E_2^{i-1,1}&\quad\text{de poids }i+1\\
E_2^{i-2,2}&\quad\text{---}\quad i+2\\[1em]
E_2^{0i}&\quad\text{---}\quad2i
\end{aligned}
\]
\margin{(lui-= muni d'une filtration \unsure{canonique} qui \unsure{correspond} aux \unsure{différents} poids\ldots)}
